\section{Methodology}
\label{sec:method}

In this section, we present the mathematical formulation of our proposed \textbf{QA-Merge} framework. We begin by outlining the low-precision quantization operator, then define the four core techniques: Quantized Centroid Calibration (QCC), Straight-Through Estimator (STE) Gating Optimization, Error-Feedback Trajectory Stabilization (EF-Smooth), and Activation Error Feedback (AEF).

\subsection{Symmetric Uniform Quantization Operator}
We model a standard symmetric uniform quantization scheme, which is highly compatible with integer-only arithmetic units in modern edge processors \cite{jacob2018quantization}. Let $x \in \mathbb{R}$ represent a continuous value (such as a weight or an activation), let $s > 0$ represent the scaling factor, and let $b$ represent the bit-width. The quantization operator $Q(x, s, b)$ maps the continuous value to a signed integer in the range $[-2^{b-1}, 2^{b-1}-1]$: $Q(x, s, b) = \text{clip}\left( \left\lfloor \frac{x}{s} \right\rceil, -2^{b-1}, 2^{b-1}-1 \right)$, where $\lfloor \cdot \rceil$ denotes the nearest-integer rounding function, and $\text{clip}(z, a, d) = \max(a, \min(z, d))$. The corresponding de-quantized value is given inline by $\tilde{x} = Q(x, s, b) \times s$.
In QA-Merge, we apply this quantization operator to all intermediate activation representations $h^{(l)}$ dynamically layer-by-layer with $b=8$ (INT8 activations). The ensembling coefficients $\boldsymbol{\alpha}$ are quantized to $b=4$ (INT4 ensembling weights) to represent 16 discrete blending levels per expert.

\textbf{Robustness to Out-of-Distribution (OOD) Activations.} Under test-time serving, incoming queries may exhibit out-of-distribution (OOD) activation profiles with ranges wildly exceeding those observed in the calibration set. If the scale factor $s_{\text{act}}$ is calibrated strictly on the absolute maximum of calibration activations, OOD outliers could induce severe register overflow or underflow. To prevent this, we employ a conservative percentile-based calibration strategy (specifically setting $s_{\text{act}}$ based on the $99.9$-th percentile of absolute activation values in the calibration set rather than the maximum) combined with the hardware-level clipping operator $\text{clip}(\cdot)$ in $Q(x, s, b)$. This clamping acts as an active hardware-level guardrail that saturates unexpected OOD outliers rather than allowing arithmetic wrap-around, guaranteeing the numerical stability of subsequent integer matrix multiplications. While percentile-based static calibration is highly effective and introduces zero runtime overhead, under extremely non-stationary OOD streams, static scale factors may still lead to excessive clipping of severe outliers. An alternative systems-level approach is \emph{dynamic activation scaling}, where the scale factor $s_{\text{act}}$ is updated on-the-fly per sample or per mini-batch (e.g., $s_{\text{act}} = \max(|x|) / 127$). While dynamic scaling completely eliminates saturation and maximizes representation range, computing the absolute maximum dynamically at each layer introduces non-trivial runtime latencies and pipeline stalls on edge microprocessors due to the need for reduction sweeps across SRAM channels. Thus, our static percentile-based calibration combined with hardware clipping strikes an optimal, low-overhead balance for standard edge serving.

\subsection{Quantized Centroid Calibration (QCC)}
In Stateless Activation-Based Latent Ensembling (\textsc{Sable}) \cite{sable2024stateless}, task-specific routing is guided by the distance between intermediate activations and reference task centroids. In high-precision floating-point, these centroids are computed as the average representation of task-specific samples. However, under extreme quantization, activations are projected onto a coarse, discrete grid. Standard float32 centroids, when rounded directly, can shift and overlap, severely eroding task separation.

To resolve this, we introduce \textbf{Quantized Centroid Calibration (QCC)}. Rather than rounding individual activations before averaging (which introduces excessive quantization noise), QCC first computes task-specific centroid means in the high-precision continuous representation space using a tiny offline calibration dataset $\mathcal{D}_k$ for each task $k \in \{1, \dots, K\}$, and then quantizes the final averaged centroid. Let $h_i^{(3)} \in \mathbb{R}^D$ represent the continuous activation extracted at Layer 3 (the frozen coordinate boundary layer) for sample $i \in \mathcal{D}_k$. The calibrated, integer-coordinate centroid $c'_k \in \mathbb{Z}^D$ is computed as: $c'_k = Q\left( \frac{1}{|\mathcal{D}_k|} \sum_{i \in \mathcal{D}_k} h_i^{(3)}, s_{act}, 8 \right)$, where $s_{act}$ is the activation scale factor. This statistically superior estimator maintains task separation and prevents centroids from merging in low precision while minimizing quantization-noise amplification. Furthermore, while standard \textsc{Sable} computes gating logits using negative squared Euclidean distance, QA-Merge utilizes scale-invariant cosine similarity in the integer coordinate space: $d_{k, b} = \frac{Q(h_b^{(3)}, s_{\text{act}}, 8) \cdot c'_k}{\| Q(h_b^{(3)}, s_{\text{act}}, 8) \|_2 \| c'_k \|_2 + \epsilon}$, where $\epsilon = 10^{-8}$ ensures numerical stability. This scale-invariant formulation prevents quantization range mismatches from corrupting the softmax boundaries.

\subsection{Quantization-Aware Gating with STE}
To route streaming queries under parameterized gating heads, we optimize routing weights $W_g \in \mathbb{R}^{K \times D}$ and biases $b_g \in \mathbb{R}^K$. Under quantization, the gating logit $z_{k, b}$ for expert $k$ under sample $b$ is computed using the quantized, INT8 representation $Q(h_b^{(3)}, s_{act}, 8)$: $z_{k, b} = \mathbf{w}_k^T Q(h_b^{(3)}, s_{act}, 8) + b_{g, k}$, where the gating weights $W_g$ and biases $b_g$ are kept in Float32 during our few-shot optimization training. However, to support a true integer-only physical processor, we apply post-training static quantization to obtain 8-bit gating weights $W_{g, \text{quant}} = Q(W_g, s_W, 8)$ and 32-bit biases $b_{g, \text{quant}} = Q(b_g, s_b, 32)$. This allows the gating layer to be computed entirely via low-power integer matrix multiplication: $z'_{k, b} = \mathbf{w}_{g, \text{quant}, k}^T Q(h_b^{(3)}, s_{\text{act}}, 8) + b_{g, \text{quant}, k}$, where $z'_{k, b}$ is computed in a 32-bit integer register. To feed these integer logits into the softmax function while preserving the original floating-point scale (since $z_{k,b} \approx z'_{k,b} \times s_z$), we scale them using the combined scale factor $s_z = s_W s_{\text{act}}$ and divide by temperature $\tau$: $\alpha_{k, b} = \exp(z'_{k, b} s_z / \tau) / \sum_{j=1}^K \exp(z'_{j, b} s_z / \tau)$.
This mathematically guarantees that the routing probability distribution is preserved identically under integer-only register computation, as detailed in Section~\ref{sec:experiments}.

\textbf{Register-Level Scale Alignment for Combined Gating.} When combining parametric routing logits $z'_{k, b}$ with the scale-invariant cosine similarity gating distances $d_{k, b} \in [-1, 1]$ prior to the softmax (i.e., $\boldsymbol{\alpha}_{\text{raw}} = \text{softmax}(\mathbf{z} + \mathbf{d})$ as shown in Appendix~\ref{app:inference_loop}), a register-level scale alignment is required because $z'_{k, b}$ is an unscaled 32-bit integer (with an implicit scale of $s_z$) while $d_{k, b}$ is a fractional value. To perform this addition entirely in low-precision integer-only registers, we project the fractional cosine similarity into a fixed-point integer representation $d'_{k, b} \in \mathbb{Z}$ by scaling it by the inverse of the logit scale factor $s_z$: $d'_{k, b} = \text{clip}\left(\left\lfloor d_{k, b} / s_z \right\rceil, -2^{15}, 2^{15}-1\right)$ inside a 16-bit signed integer register. This aligns both terms to a unified scale of $s_z$, allowing the addition to be executed via a single-cycle integer instruction: $z''_{k, b} = z'_{k, b} + d'_{k, b}$. The aligned integer sum is then fed into the softmax by scaling the exponent by $s_z/\tau$: $\alpha_{k, b} = \exp(z''_{k, b} s_z / \tau) / \sum_{j=1}^K \exp(z''_{j, b} s_z / \tau)$. This protocol avoids any runtime floating-point conversions, preserving our fully integer-only pipeline.

During offline few-shot optimization, we update $W_g$ and $b_g$ by minimizing the downstream cross-entropy loss. However, because the rounding function $\lfloor \cdot \rceil$ has zero derivative almost everywhere, standard backpropagation fails due to gradient vanishing.

To enable gradient-based training through these discrete rounding boundaries, we employ the \textbf{Straight-Through Estimator (STE)} \cite{bengio2013estimating}. In the backward pass, the gradient of the rounding operator is approximated as the identity mapping $\frac{\partial \tilde{h}}{\partial h} \approx 1$.
This allows gradients to flow natively back into the gating weights, preventing dead routing nodes and allowing the gating head to learn optimal separation boundaries directly adapted to the integer activation landscape.

\subsection{Error-Feedback Trajectory Stabilization (EF-Smooth)}
In deep coordinate-space architectures, blending weights $\boldsymbol{\alpha}$ are computed at each layer to linearly interpolate representations from specialized experts. Under low-precision constraints, we quantize these coefficients to 4-bit precision (INT4), restricting them to 16 discrete levels of expert contributions. Rounding these coefficients at each layer introduces significant discretization noise. Across a 14-layer deep network, these rounding errors accumulate, leading to representational drift and trajectory jitter.

To stabilize these ensembling trajectories, we propose \textbf{Error-Feedback Trajectory Stabilization (EF-Smooth)}, inspired by error-diffusion algorithms in digital signal processing. EF-Smooth tracks the rounding error of ensembling coefficients at layer $l$ and injects it back into the routing weight of layer $l+1$ as a correction.

Standard simplex projections scale ensembling weights off the discrete 4-bit grid, violating low-precision hardware constraints. To solve this, we employ a \textbf{Discrete Simplex Projection} algorithm (similar to Hamilton's method of apportionment), which projects continuous weights onto strictly discrete 4-bit integers $q_k \in \{0, \dots, 15\}$ summing exactly to $15$: $\sum_{k=1}^K q_k = 15$. The final 4-bit quantized coefficient is then represented as $\tilde{\alpha}_{k}^{(l)} = q_k / 15$. 

Specifically, the \text{DiscreteSimplexProj} operator maps continuous weights $\boldsymbol{\alpha}$ and a target bit-width $b$ (with $L = 2^b - 1 = 15$ levels) onto the discrete simplex. First, the vector is projected onto the standard continuous simplex: $\alpha'_k = \max(0, \alpha_k) / \sum_{j=1}^K \max(0, \alpha_j)$. Second, the normalized values are scaled and truncated to obtain initial allocations $q_k = \lfloor \alpha'_k L \rfloor$ and fractional remainders $r_k = \alpha'_k L - q_k$. Third, the shortfall $S = L - \sum_{k=1}^K q_k$ is calculated. In standard Hamilton apportionment, allocations are adjusted by sorting expert indices $\{1, \dots, K\}$ in descending order of remainders $r_k$ and incrementing $q_k \leftarrow q_k + 1$ for the $S$ experts with the largest remainders. However, sorting expert remainders has $O(K \log K)$ complexity, which is highly inefficient on specialized edge vector hardware due to pipeline branches. 

To bypass sorting while preserving permutation invariance and remainder-magnitude sensitivity, we introduce \textbf{Permutation-Invariant Single-Pass Apportionment (PI-SPA)}. First, we perturb remainders $r_k$ with a tiny static unique expert identifier: $r'_k = r_k + \epsilon \cdot \text{ID}_k$, where $\text{ID}_k \in \{0, \dots, K-1\}$ is the static expert index and $\epsilon \ll 1/L$ (typically $\epsilon = 10^{-5}$). Since each expert has a unique static ID independent of its memory listing order, the perturbed remainders $r'_k$ are strictly unique.

Second, we determine the $S$-th largest perturbed remainder as our allocation threshold: $\theta = \text{Select}\left(\{r'_j\}_{j=1}^K, S\right)$ in $O(K)$ time (or $\theta = \infty$ if $S=0$). Since $K$ is small on edge hardware (e.g., $K \le 8$), this is executed efficiently using parallel Selection or tiny branchless bitonic networks on specialized vector pipelines, avoiding pipeline branches.

Third, allocations are incremented for the $S$ experts meeting this threshold: $q_k \leftarrow q_k + 1$ for all $k$ where $r'_k \geq \theta$. Strict uniqueness of $r'_k$ guarantees that exactly $S$ experts meet this threshold, ensuring the shortfall is fully allocated and the ensembling weights $\tilde{\alpha}_k = q_k / L$ strictly sum to 1.0. 

Since $\text{ID}_k$ is static, the selected experts are strictly invariant to input expert list permutations, eliminating compilation fragility. Since $\epsilon \ll 1/L$, ordering is dominated by remainders $r_k$; static IDs only resolve ties, preserving remainder-magnitude sensitivity. While static tie-breaking introduces a negligible deterministic statistical bias under exact ties, exact ties are exceptionally rare and have zero observed impact on accuracy. 

Let $\boldsymbol{\alpha}_b^{(l)} \in [0, 1]^K$ represent the continuous routing weight vector computed at layer $l$. The quantization error vector $\mathbf{e}_{\alpha, b}^{(l)}$ at layer $l$ is defined as $\mathbf{e}_{\alpha, b}^{(l)} = \boldsymbol{\alpha}_{b, \text{corrected}}^{(l)} - \tilde{\boldsymbol{\alpha}}_b^{(l)}$. For the subsequent layer $l+1$, we add the error feedback discounted by a decay factor $\beta \in [0, 1]$ to form the corrected continuous routing vector: $\boldsymbol{\alpha}_{b, \text{corrected}}^{(l+1)} = \boldsymbol{\alpha}_b^{(l+1)} + \beta \mathbf{e}_{\alpha, b}^{(l)}$. The final discrete, quantized ensembling coefficient vector $\tilde{\boldsymbol{\alpha}}_b^{(l+1)}$ is then obtained using the discrete simplex projection operator: $\tilde{\boldsymbol{\alpha}}_b^{(l+1)} = \text{DiscreteSimplexProj}\left( \boldsymbol{\alpha}_{b, \text{corrected}}^{(l+1)}, 4 \right)$.
This error diffusion mechanism keeps the average representational blending path extremely close to the optimal continuous trajectory, effectively neutralizing discretization noise with near-zero computational cost.

\subsection{Activation Error Feedback (AEF)}
In addition to weight trajectory smoothing, a critical challenge in low-precision representation ensembling is the \emph{Small-Step Quantization Bottleneck}. Because the adapter scaling factor $\gamma_V = 0.05$ is very small, the dynamic representation update at each layer (the ``pull'' vector) is extremely small ($\approx 0.01$). Under standard INT8 activation quantization (where the grid spacing is $\approx 0.03$), these updates are frequently rounded to zero: $Q(\tilde{h}^{(l-1)} + \text{pull}, s_{\text{act}}, 8) = Q(\tilde{h}^{(l-1)}, s_{\text{act}}, 8)$.
To prevent this information collapse, we introduce \textbf{Activation Error Feedback (AEF)}. AEF tracks the sub-grid quantization error of representations layer-by-layer and adds it back as a residually accumulated bias to the next layer's update:
\begin{equation}\label{eq:aef_unquantized}
    h^{(l)}_{\text{unquantized}} = \tilde{h}^{(l-1)} + \text{pull} + e^{(l-1)}_{\text{act}}
\end{equation}
\begin{equation}\label{eq:aef_quantized}
    \tilde{h}^{(l)} = Q(h^{(l)}_{\text{unquantized}}, s_{\text{act}}, 8) \times s_{\text{act}}
\end{equation}
\begin{equation}\label{eq:aef_error}
    e^{(l)}_{\text{act}} = h^{(l)}_{\text{unquantized}} - \tilde{h}^{(l)}
\end{equation}
This closed-loop accumulation guarantees that tiny, sample-specific updates are not discarded; instead, they are accumulated residually until they cross the quantization grid threshold, triggering a discrete step in the integer space.\footnote{While the formulation in Eqs.~\ref{eq:aef_unquantized}--\ref{eq:aef_error} assumes a constant activation scale factor $s_{\text{act}}$ across layers for ease of exposition, in practical hardware deployments, dynamic scale factors $s_{\text{act}}^{(l)}$ are often adjusted layer-by-layer to maximize the dynamic range of INT8 activations. In such scenarios, the accumulated error buffer $\mathbf{e}_{\text{act}}^{(l-1)}$ must be dynamically rescaled to the current layer's scale factor: $\mathbf{e}_{\text{act}, \text{realigned}}^{(l-1)} = \mathbf{e}_{\text{act}}^{(l-1)} \times \left( s_{\text{act}}^{(l-1)} / s_{\text{act}}^{(l)} \right)$, which is efficiently implemented in hardware using fixed-point dyadic multiplication (integer shift-and-scale) as detailed in Appendix~\ref{app:hardware_systems}.} 

To demonstrate the theoretical rigor of AEF, we provide a formal mathematical proof of its error-bounding properties. We show that AEF has a telescoping property that bounds the cumulative activation quantization error relative to the quantized-pull accumulated trajectory (defined using the intermediate quantized states).
\begin{theorem}[Telescoping Bounded Representational Error of AEF]\label{thm:aef_error_bound}
Let $h^{(l)}_{\text{unquantized}}$ be propagated via AEF as defined in Eqs.~\ref{eq:aef_unquantized}--\ref{eq:aef_error}. Under a frozen boundary activation $\tilde{h}^{(3)}$, the accumulated representation at the final layer $L$ satisfies the following relation relative to the quantized-pull accumulated trajectory $\tilde{h}^{(3)} + \sum_{l=4}^L \text{pull}^{(l)}$ (which we note is a local quantized-state trajectory defined using intermediate quantized states, and may diverge from the true continuous unquantized trajectory $h^{(3)}_{\text{float}} + \sum_{l=4}^L \text{pull}^{(l)}_{\text{float}}$ due to feedback-driven pull calculations based on quantized intermediate states):
\begin{equation}
    \tilde{h}^{(L)} = \tilde{h}^{(3)} + \sum_{l=4}^L \text{pull}^{(l)} - e^{(L)}_{\text{act}}
\end{equation}
Furthermore, assuming uniform rounding bounds where each coordinate of the error vector satisfies $|e_{\text{act}, d}^{(l)}| \leq \frac{s_{\text{act}}}{2}$, the total cumulative representational error relative to this quantized-pull trajectory at any depth $L$ is strictly bounded by:
\begin{equation}
    \left\| \tilde{h}^{(L)} - \left( \tilde{h}^{(3)} + \sum_{l=4}^L \text{pull}^{(l)} \right) \right\|_2 \leq \frac{s_{\text{act}} \sqrt{D}}{2}
\end{equation}
\end{theorem}
\begin{proof}
See Appendix~\ref{app:aef_proof} for the complete mathematical derivation.
\end{proof}
This formalizes why AEF stabilizes representational trajectories: by acting as a perfect closed-loop error integrator, the cumulative error is bounded solely by the single-layer quantization step of the final boundary, completely eliminating error accumulation and drift across arbitrary network depths. 

We discuss the subtle distinction between the bounded numerical error of AEF and continuous trajectory divergence in Appendix~\ref{app:aef_proof}. Empirically, we show in Section~\ref{sec:experiments} that this divergence remains highly benign.

